\documentclass[10pt,conference]{IEEEtran}

\usepackage[utf8]{inputenc}
\usepackage[T1]{fontenc}
\usepackage{amsmath,amssymb,amsfonts}
\usepackage{algorithmic}
\usepackage{graphicx}
\usepackage{textcomp}
\usepackage{xcolor}
\usepackage{booktabs}
\usepackage{tabularx}
\usepackage{url}
\usepackage{cite}
\usepackage{balance}

\hyphenation{re-chunk-ing spatio-tem-po-ral par-ti-tion-ing me-ga-drought}

\begin{document}

\title{Distributed Spatio-Temporal Processing of Massive Climate Grids for Fire Weather Index Estimation: State-of-the-Art and Architecture Design}

\author{
    \IEEEauthorblockN{Agust\'in Pizarro, Eduardo Valenzuela, Benjam\'in Z\'u\~niga}
    \IEEEauthorblockA{\textit{School of Informatics and Telecommunications} \\
    \textit{Universidad Diego Portales (UDP)}\\
    Santiago, Chile \\
    Email: \{agustin.pizarro, eduardo.valenzuela, benjamin.zuniga\}@mail.udp.cl}
    \and
    \IEEEauthorblockN{Nicol\'as Hidalgo (Instructor)}
    \IEEEauthorblockA{\textit{School of Informatics and Telecommunications} \\
    \textit{Universidad Diego Portales (UDP)}\\
    Santiago, Chile \\
    Email: nicolas.hidalgoc@mail.udp.cl}
}

\maketitle

\begin{abstract}
Wildfire danger in Chile has reached catastrophic levels due to a 15-year uninterrupted megadrought and extreme synoptic conditions adhering to the ``30-30-30'' rule. Operational fire danger assessment relies on the Canadian Forest Fire Weather Index (FWI) System, which is historically evaluated exclusively at sparse, physically vulnerable automatic weather stations (AWS), leaving hazardous spatial blind spots across complex Andean terrain. While high-resolution gridded climate products such as CR2MET v2.5 ($\sim$5~km) provide wall-to-wall spatial coverage, audit of official archives reveals they supply only daily precipitation and temperature extremes, entirely lacking relative humidity and surface wind. To achieve high-resolution, decadal FWI tracking across Chile, a multi-source distributed data pipeline is required to couple CR2MET with hourly ERA5-Land reanalysis and the 35-model NASA NEX-GDDP-CMIP6 climate projection ensemble.

Crucially, while the arithmetic evaluation of the recursive FWI formulas requires negligible CPU time per cell, the system-level challenge constitutes an authentic distributed systems bottleneck. We formalize this challenge through three distributed dimensions: (i) the \textit{Distributed Corner-Turn / Rechunking Shuffle Problem}, wherein source data stored in time-sliced spatial rasters (\texttt{[time, lat, lon]}) must be transposed across a cluster into cell-major time-series (\texttt{[cell\_id, time]}) to satisfy a 44-year non-commutative, stateful temporal recursion; (ii) multi-resolution spatio-temporal alignment across heterogeneous grids ($0.05^\circ$, $0.1^\circ$, and $0.25^\circ$); and (iii) long-term Markovian state persistence and checkpointing. 

In this paper, we conduct a rigorous, verified state-of-the-art review covering 18 foundational and contemporary works across distributed processing engines, scientific array frameworks, and fire climatology. We provide a taxonomic comparison demonstrating that existing array systems fail to address multi-decadal, time-recursive stateful scans. Finally, we design a 4-stage distributed architecture using Apache Spark over columnar Apache Parquet with zero-shuffle spatial-first partitioning, evaluated against a Dask/Zarr baseline in a resource-constrained containerized Docker cluster.
\end{abstract}

\begin{IEEEkeywords}
Distributed Computing, Big Data Systems, Apache Spark, Columnar Storage, Fire Weather Index, Climate Tensors, Rechunking Shuffle, CR2MET, ERA5-Land.
\end{IEEEkeywords}

\section{Introduction and Problem Formalization}

\subsection{Context and Motivation}
During the last decade, Central and Southern Chile have faced devastating wildfire megaseasons, most notably the 2017 ``Tormenta de Fuego'' ($>570,000$~ha burned) and the 2023 disaster ($>430,000$~ha, 26 fatalities)~\cite{cordero2024extreme, bowman2019human}. These catastrophes stem from the interplay of a 15-year uninterrupted megadrought~\cite{garreaud2020central}, extensive contiguous monoculture plantations of flammable pine and eucalyptus, and recurrent heatwaves governed by the empirical ``30-30-30'' rule (temperature $T > 30^\circ\text{C}$, relative humidity $RH < 30\%$, and 10~m wind speed $W > 30\text{ km/h}$).

Operational fire management agencies, including the Chilean National Forestry Corporation (CONAF) and the National Disaster Prevention and Response Service (SENAPRED), rely on the Canadian Forest Fire Danger Rating System (CFFDRS) and its core subsystem, the Fire Weather Index (FWI)~\cite{vanwagner1987development}. The FWI models the physical drying and moisture dynamics across three distinct forest floor organic layers:
\begin{enumerate}
    \item \textbf{Fine Fuel Moisture Code (FFMC):} Surface litter and cured fine fuels (fast response: 16--24~h timelag).
    \item \textbf{Duff Moisture Code (DMC):} Moderately compacted decomposing organic material ($\sim$12-day timelag).
    \item \textbf{Drought Code (DC):} Deep, compacted organic soil layers, serving as a seasonal indicator of cumulative megadrought ($\sim$52-day timelag).
\end{enumerate}
These fuel moisture codes drive three intermediate fire behavior indices: the Initial Spread Index (ISI), Buildup Index (BUI), and the composite Fire Weather Index (FWI) representing frontal flame intensity~\cite{wotton2009interpreting}.

\subsection{The Spatial Monitoring Bottleneck and CR2MET Shortfall}
Historically, FWI is evaluated exclusively at physical Automated Weather Stations (AWS) operated by the Chilean Weather Service (DMC) and the Water Directorate (DGA). However, Chile's rugged Andean topography creates steep microclimatic gradients across narrow longitudinal spans. The AWS network is sparse, predominantly urban or coastal, and subject to sensor destruction during wildfires. Consequently, critical rural forest interfaces remain unmonitored ``blind spots'' where ignitions propagate undetected.

To overcome AWS sparsity, climate researchers developed gridded products such as \textbf{CR2MET v2.5}~\cite{boisier2023cr2met}, which models daily climate over continental Chile on a regular $0.05^\circ \times 0.05^\circ$ grid ($\sim$5~km resolution, comprising $\sim$20,000 active land cells from 1960 to 2021). 

However, an audit of the official CR2MET v2.5 distribution (Zenodo record \texttt{10.5281/zenodo.7529682}) reveals a critical operational shortfall: **CR2MET v2.5 exclusively provides daily precipitation (\texttt{pr}) and extreme temperatures (\texttt{tmin}, \texttt{tmax})**. It does \textit{not} provide relative humidity ($RH$) or wind speed ($W$), which are mandatory terms in the differential drying equations of FFMC and ISI. 

Therefore, constructing a continuous wall-to-wall FWI database over Chile is impossible using CR2MET alone. It strictly requires a \textbf{multi-source distributed data fusion pipeline} that couples CR2MET's $0.05^\circ$ topographic precipitation/temperature with hourly atmospheric moisture ($T_d$) and vector wind components ($u_{10}, v_{10}$) from ECMWF's **ERA5-Land** reanalysis ($0.1^\circ$ grid, 1950--present)~\cite{munoz2021era5land}, and projects historical risk forward across the 35 General Circulation Models (GCMs) and 4 Shared Socioeconomic Pathways (SSPs) of the **NASA NEX-GDDP-CMIP6** ensemble ($0.25^\circ$ grid, 1950--2100)~\cite{quilcaille2023fire}.

\subsection{The Distributed Systems Challenge: Why is FWI Genuine Big Data?}
At first glance, evaluating the empirical equations of the Canadian FWI might appear computationally trivial. For a single grid cell on a single day, the algebraic formulas for FFMC, DMC, and DC require fewer than 60 floating-point operations (FLOPs). Consequently, evaluating $3.2 \times 10^8$ cell-days on a modern CPU executing at several GFLOPs would theoretically take only a few seconds if all data resided in contiguous L1/L2 cache. 

The true challenge is fundamentally an **I/O, memory hierarchy, and distributed network communication bottleneck**, characterized by three core system properties:

\subsubsection{The Distributed Corner-Turn (Rechunking Shuffle) Problem}
Let $\mathcal{S}_{c,t} = \langle FFMC_{c,t}, DMC_{c,t}, DC_{c,t} \rangle \in \mathbb{R}^3$ denote the state vector of cell $c \in \{1, \dots, N_c\}$ at day $t \in \{1, \dots, T\}$, and let $\mathbf{x}_{c,t} = \langle T_{c,t}, RH_{c,t}, W_{c,t}, P_{c,t} \rangle \in \mathbb{R}^4$ be the input forcing. The state transition is governed by:
\begin{equation}
\label{eq:recurrence}
\mathcal{S}_{c,t} = \Phi\left(\mathcal{S}_{c,t-1},\, \mathbf{x}_{c,t},\, \text{month}(t),\, \text{lat}(c)\right)
\end{equation}
Because the physical drying functions inside $\Phi$ are exponential, piecewise-linear, and conditioned on rainfall saturation thresholds, $\Phi$ is \textbf{strictly non-commutative and non-associative across time}:
\begin{equation}
\Phi(\Phi(\mathcal{S}, \mathbf{x}_1), \mathbf{x}_2) \neq \Phi(\Phi(\mathcal{S}, \mathbf{x}_2), \mathbf{x}_1)
\end{equation}
Thus, time cannot be folded via a parallel tree reduction ($O(\log T)$). The temporal trajectory must be evaluated sequentially as a prefix scan of length $T$ for each spatial cell $c$.

Crucially, raw climate datasets (NetCDF-4/GRIB2) are archived in \textbf{time-major order} (\texttt{[time, lat, lon]}), where each 2D file represents a spatial map at time $t$. If a single node attempts to read the 44-year time series of cell $c$, it must perform $T \approx 16,000$ non-contiguous, fragmented disk seeks across thousands of compressed NetCDF chunk boundaries. This produces a \textit{Read Amplification Factor} exceeding $10^3\times$, collapsing throughput. 

Conversely, if $P$ distributed worker nodes ingest time-sliced blocks in parallel, no single worker possesses the contiguous history of any spatial cell. Grouping observations by cell requires an **all-to-all distributed transpose (corner-turn shuffle)** that transfers:
\begin{equation}
V_{\text{shuffle}} \approx \left(1 - \frac{1}{P}\right) \cdot N_c \cdot T \cdot |\mathbf{x}|
\end{equation}
across cluster interconnects. For hundreds of gigabytes of climate tensors, an unoptimized shuffle triggers severe network saturation, JVM memory pressure, and out-of-memory (OOM) failures.

\subsubsection{Multi-Resolution Spatio-Temporal Join}
Coupling CR2MET ($0.05^\circ$, daily), ERA5-Land ($0.1^\circ$, hourly), and CMIP6 ($0.25^\circ$, multi-model daily) requires continuous bilinear regridding and psychrometric derivation of relative humidity from dewpoint temperature ($T_d$) via the Magnus-Tetens formula. This spatial join expands the evaluation domain to $>7.7 \times 10^9$ hourly state updates for historical reanalysis and $>1.5 \times 10^{10}$ projections across the 35 CMIP6 GCMs.

\subsubsection{Stateful Checkpointing and Fault Tolerance}
Because the Drought Code has a 52-day memory and accumulates deficits across multi-year droughts, worker failures during long-running batch iterations could cause catastrophic recomputation of RDD lineages. The engine must decouple long time-series into bounded epochs with distributed state checkpointing.

\section{Related Work}

We classify the state-of-the-art into three distinct research pillars, totaling 18 verified references evaluated according to the criteria defined in~\cite{paperinstructions}.

\subsection{Distributed Big Data Systems and Columnar Formats}
The foundation of modern distributed data processing originated with Dean and Ghemawat's **MapReduce**~\cite{dean2008mapreduce}. While MapReduce introduced resilient large-scale parallel primitives, its mandatory materialization of intermediate states to disk after every map and reduce phase renders multi-decadal recursive scans computationally intractable. 

To overcome disk I/O bottlenecks, Zaharia et al. developed **Apache Spark** and Resilient Distributed Datasets (RDDs)~\cite{zaharia2012resilient}, introducing in-memory fault-tolerant abstractions. Spark enables fine-grained control over partitioning through custom hash and range partitioners, allowing narrow lineage dependencies that avoid network shuffles. 

Building upon RDDs, Armbrust et al. introduced **Spark SQL**~\cite{armbrust2015sparksql}, which incorporates the Catalyst query optimizer and the Tungsten execution engine. Catalyst automates logical-to-physical query plan transformations, pushdown predicates, and columnar projection, significantly reducing serialization overhead between the JVM and native memory. For streaming workloads, Armbrust et al. presented **Structured Streaming**~\cite{armbrust2018structured}, providing declarative stateful APIs (\texttt{mapGroupsWithState}) with end-to-end exactly-once guarantees. 

In parallel, Carbone et al. investigated state management within **Apache Flink**~\cite{carbone2017state}, demonstrating that continuous event-driven stream processing with asynchronous Chandy-Lamport state snapshotting can achieve sub-second stateful updates. 

Complementing execution engines, Melnik et al. developed **Dremel**~\cite{melnik2010dremel}, introducing nested columnar storage schemes that served as the architectural blueprint for **Apache Parquet**. Parquet utilizes dictionary encoding, bit-packing, and run-length encoding (RLE) to deliver high compression ratios alongside embedded min/max block statistics for efficient data skipping (*predicate pushdown*).

\subsection{Scientific Array Computing and Geospatial Engines}
The adaptation of big data frameworks to multi-dimensional scientific arrays has been widely explored. Palamuttam et al. introduced **SciSpark**~\cite{palamuttam2015scispark}, extending Spark with a specialized array abstraction (\texttt{sciRDD}) to perform distributed weather event detection from NetCDF files. However, SciSpark's partitioner divides data across spatial 2D slices at fixed timestamps, which fails to support longitudinal, time-recursive scans without continuous network re-partitioning. 

Hu et al. developed **ClimateSpark**~\cite{hu2018climatespark}, an in-memory distributed computing framework tailored to big climate data analytics. ClimateSpark integrates a spatiotemporal index over HDFS blocks to optimize reading NetCDF chunks. Similarly, Li et al.~\cite{li2016spatiotemporal} demonstrated spatiotemporal indexing structures on MapReduce to reduce read amplification in MERRA reanalysis datasets. Nonetheless, both frameworks target commutative and associative spatio-temporal reductions (such as global monthly means or regional variance) rather than non-commutative stateful prefix scans.

For vector and raster geospatial operations, Yu, Zhang, and Sarwat introduced **GeoSpark** (now **Apache Sedona**)~\cite{yu2018spatial}, which embeds spatial indexing (R-tree and Quad-tree) and spatial partitioning into Spark. GeoSpark achieves scalable spatial joins between point sets and polygons, which is critical for intersecting gridded climate rasters with administrative commune boundaries and active fire detections. 

On the public cloud, Gorelick et al. established **Google Earth Engine (GEE)**~\cite{gorelick2017google}, a planetary-scale platform for earth observation data. While GEE provides seamless spatial map algebra over multi-petabyte satellite archives, its computational model is fundamentally optimized for stateless, per-pixel 2D image operations. Expressing multi-decadal recursive scans in GEE via \texttt{ee.ImageCollection.iterate()} rapidly exceeds platform recursion limits and user memory budgets, causing pipeline aborts.

Within the scientific Python ecosystem, Rocklin presented **Dask**~\cite{rocklin2015dask}, which, combined with Hoyer and Hamman's **xarray**~\cite{hoyer2017xarray}, forms the foundation of the **Pangeo** initiative. Abernathey et al.~\cite{abernathey2021cloud} synthesized the requirements for cloud-native scientific repositories, identifying the fundamental tension between time-chunked and space-chunked layouts in **Zarr** and NetCDF storage. They developed *Rechunker*, a distributed library that coordinates intermediate, two-phase array transpositions in object stores.

\subsection{Fire Weather Climatology and Chile Wildfire Dynamics}
The physical underpinnings of the Canadian FWI system were formalized by Van Wagner~\cite{vanwagner1987development}. Recently, the Canadian Forest Service Fire Danger Group released the **FWI2025** update~\cite{cfs2025fwi}, which introduces the *Hourly Fine Fuel Moisture Code* (Hourly FFMC) to capture diurnal drying cycles, overnight moisture recovery, and shifting peak burning windows. 

At the global scale, Vitolo et al.~\cite{vitolo2020era5} established the Global ECMWF Fire Forecasting (GEFF) reanalysis dataset using ERA5 at a $0.25^\circ$ ($\sim$28~km) spatial resolution. Similarly, Quilcaille et al.~\cite{quilcaille2023fire} computed global FWI projections under CMIP6 scenarios from 1850 to 2100 at coarse grid scales ($1^\circ$--$2^\circ$). 

In Chile, Garreaud et al.~\cite{garreaud2020central} characterized the dynamical atmospheric drivers of the Central Chile megadrought, demonstrating an ongoing $30\%$ regional rainfall deficit. Boisier~\cite{boisier2023cr2met} produced the CR2MET v2.5 high-resolution precipitation and temperature grids. Cordero et al.~\cite{cordero2024extreme} and Bowman et al.~\cite{bowman2019human} proved that extreme fire weather driven by ENSO and anthropogenic warming tripled the burned area in Central Chile over the last decade, underscoring the urgent necessity of high-resolution, continuous fire hazard tracking.

\begin{table*}[t]
\centering
\caption{Taxonomic Comparison of Distributed Processing Systems on Multi-Decadal Gridded Climate Workloads}
\label{tab:taxonomy}
\small
\begin{tabularx}{\textwidth}{l p{2.0cm} p{2.0cm} X X p{3.8cm}}
\toprule
\textbf{System / Framework} & \textbf{Data Model} & \textbf{Primary Storage Layout} & \textbf{Stateful Temporal Recursion} & \textbf{Shuffle Overhead on Corner-Turn} & \textbf{Primary Research Gap vs. Our Approach} \\
\midrule
\textbf{MapReduce}~\cite{dean2008mapreduce} & Key-Value pairs & Flat HDFS files (row-oriented) & No native support (disk-bound loops) & Extreme (disk-spill on all map-reduce barriers) & Prohibitive I/O latency when chaining $>16,000$ sequential daily timesteps. \\
\addlinespace
\textbf{SciSpark}~\cite{palamuttam2015scispark} & \texttt{sciRDD} multi-dim arrays & NetCDF on HDFS (spatial 2D slices) & Not supported (stateless transformations) & High (requires repartitioning across time) & Partitions by timestamp; cannot execute time-continuous cell scans without global shuffle. \\
\addlinespace
\textbf{ClimateSpark}~\cite{hu2018climatespark} & Spatiotemporal indexed RDDs & Chunked NetCDF on HDFS & Associative reductions only (means/variances) & Moderate (optimized for spatiotemporal box queries) & Designed for aggregation queries; lacks stateful iterator primitives for recursive moisture accumulation. \\
\addlinespace
\textbf{Apache Sedona}~\cite{yu2018spatial} & Spatial RDD / Geometry & Vector GeoJSON, Shapefiles, WKT & Stateless vector spatial joins & Low for vector joins; unadapted for raw 3D arrays & Excellent for polygon/point joins (FIRMS/CONAF), but not designed for 3D tensor time-series rechunking. \\
\addlinespace
\textbf{Google Earth Engine}~\cite{gorelick2017google} & Planetary-scale Image Collections & Cloud-optimized 2D tiles (spatial-major) & Severely constrained (\texttt{iterate} depth limits) & Hidden (server-side tile computation) & Cannot execute 44-year recursive loops due to memory quotas and iteration recursion stack limits. \\
\addlinespace
\textbf{Dask + Xarray}~\cite{rocklin2015dask, abernathey2021cloud} & N-dimensional blocked arrays & Zarr / NetCDF on Object Storage (S3) & Manual (via \texttt{map\_blocks} over pre-chunked time) & High without specialized two-phase Rechunker & Requires explicit external rechunking pipeline; lacks built-in query optimizer and streaming state management. \\
\addlinespace
\textbf{GEFF / CEMS Fire}~\cite{vitolo2020era5} & Fortran monolithic binary pipeline & Coarse GRIB/NetCDF ($0.25^\circ$) & Sequential in-memory execution per gridpoint & None (runs on single HPC batch nodes) & Monolithic, non-distributed pipeline; spatial resolution ($0.25^\circ \approx 28$~km) misses Chilean mountain valleys. \\
\addlinespace
\textbf{Our Proposed Architecture} & \textbf{Spatial-First Partitioned DataFrames} & \textbf{Columnar Apache Parquet (ZSTD)} & \textbf{Native stateful prefix scan (Arrow + Numba JIT)} & \textbf{Zero-shuffle during scan phase} (single one-time repartition) & \textbf{Fills the gap:} Combines multi-source spatial fusion, one-time corner-turn, zero-shuffle stateful scans, and streaming updates. \\
\bottomrule
\end{tabularx}
\end{table*}

\section{Discussion of Related Work and Research Gap}

Table~\ref{tab:taxonomy} provides a systematic taxonomic comparison across candidate distributed processing frameworks. 

An inspection of the literature reveals a clear, unaddressed \textbf{Research Gap at the intersection of Big Data Systems and Physical Environmental Modeling}:
\begin{enumerate}
    \item \textbf{The Reductive Bias of Big Climate Systems:} Frameworks such as ClimateSpark~\cite{hu2018climatespark}, SciHadoop~\cite{li2016spatiotemporal}, and standard Spark SQL pipelines are engineered under the assumption that climate analytics consist of \textit{associative and commutative aggregations} (e.g., temporal averages, regional bounding box sums, or climate anomaly indexing). These operations scale seamlessly via tree reductions ($O(\log T)$). In sharp contrast, fire risk modeling requires an inherently non-commutative, non-associative prefix recurrence (Eq.~\ref{eq:recurrence}) that forces a sequential temporal scan over decades.
    \item \textbf{The Storage Layout / Computation Mismatch:} Existing global datasets (such as GEFF~\cite{vitolo2020era5} and CMIP6-FWI~\cite{quilcaille2023fire}) circumvent distributed systems by running coarse monolithic Fortran/C codes on high-memory HPC nodes. Conversely, cloud-native earth observation engines (like Google Earth Engine~\cite{gorelick2017google}) optimize for 2D spatial map tiles at individual timestamps. When scaled to high spatial resolution ($0.05^\circ$) over 44--75 years ($1.6 \times 10^4$ to $6.5 \times 10^5$ timesteps), neither paradigm scales: HPC nodes run out of memory, while GEE exhausts recursion execution stacks.
    \item \textbf{The Multi-Source Fusion Void:} No existing system reconciles high-resolution topographic precipitation/temperature grids with hourly atmospheric reanalysis and global multi-model projection ensembles in a unified distributed execution graph.
\end{enumerate}

Our work directly targets this gap by formulating a \textbf{Two-Phase Architecture: One-Time Distributed Corner-Turn followed by Zero-Shuffle Vectorized Stateful Processing}.

\section{Proposed Solution Architecture}

\subsection{Data Ecosystem and Volume Specifications}
The architecture integrates six verified heterogeneous data sources, summarized in Table~\ref{tab:datasources}. The aggregated historical and projection volume spans hundreds of gigabytes of compressed tensors, expanding to $>2$~TB uncompressed, validating the project as a genuine Big Data processing scenario.

\begin{table}[t]
\centering
\caption{Integrated Data Ecosystem and Storage Sizing}
\label{tab:datasources}
\footnotesize
\begin{tabularx}{\columnwidth}{l X c r}
\toprule
\textbf{Dataset} & \textbf{Variables} & \textbf{Grid / Freq.} & \textbf{Volume} \\
\midrule
CR2MET v2.5 & \texttt{pr, tmin, tmax} & $0.05^\circ$ / Day & 9.3 GB \\
ERA5-Land & \texttt{t2m, d2m, u, v, tp} & $0.10^\circ$ / Hour & 165 GB \\
NEX-GDDP & \texttt{tas, pr, hurs, sfcW} & $0.25^\circ$ / CMIP6 & 420 GB \\
CEMS FWI & Historical \texttt{FWI, FFMC, DC} & $0.25^\circ$ / Day & 12 GB \\
NASA FIRMS & Active fire (\texttt{frp}, pts) & 375 m / Daily & 450 MB \\
CONAF Cat. & Ignition logs \& area & Point / Event & 85 MB \\
\bottomrule
\end{tabularx}
\end{table}

\subsection{Four-Stage Distributed Workflow}
The workflow is divided into four modular processing stages:

\subsubsection{Stage 1: Multi-Source Distributed Ingestion and Fusion}
Worker nodes read spatial partitions of CR2MET NetCDF files alongside ERA5-Land reanalysis from object storage. For each hourly timestamp in ERA5-Land, saturation vapor pressure ($e_s$) and actual vapor pressure ($e$) are computed from $T$ and $T_d$ using the Magnus-Tetens relationship:
\begin{equation}
e(T) = 6.112 \cdot \exp\left(\frac{17.67 \cdot T}{T + 243.5}\right) \quad [\text{hPa}]
\end{equation}
\begin{equation}
RH = \min\left(100.0,\, 100.0 \cdot \frac{e(T_d)}{e(T)}\right) \quad [\%]
\end{equation}
Wind speed magnitude is derived as $W = 3.6 \sqrt{u_{10}^2 + v_{10}^2}\text{ km/h}$. Using Apache Spark DataFrames, hourly values are bilinearly interpolated onto CR2MET's $0.05^\circ$ reference grid and assigned a spatial key based on the \textbf{Uber H3 Hierarchical Hexagonal Spatial Index} (resolution 6, $\sim$3.2~km aperture).

\subsubsection{Stage 2: Distributed Corner-Turn and Columnar Materialization}
To permanently resolve the time-major read amplification bottleneck, Spark executes a controlled global partition:
\begin{equation*}
\begin{split}
\texttt{df.repartitionByRange}(P,\, \texttt{"h3\_id"}) \\
\texttt{.sortWithinPartitions("h3\_id", "time")}
\end{split}
\end{equation*}
where $P$ is the cluster parallelism factor. The transposed data is materialized into \textbf{Apache Parquet files compressed with Zstandard (ZSTD)}, partitioned by \texttt{h3\_id}. Each Parquet file encapsulates the entire contiguous multi-decadal time-series for a local spatial cluster of cells. Embedded dictionary encoding and min/max block statistics eliminate redundant metadata reads.

\subsubsection{Stage 3: Zero-Shuffle Vectorized Stateful Scan}
Because each executor now possesses the complete temporal trajectory for its assigned spatial partition, FWI evaluation proceeds with \textbf{zero inter-node network communication}:
\begin{equation*}
\begin{split}
\texttt{df.groupBy("h3\_id")} \\
\texttt{.applyInPandas(compute\_fwi\_numba, schema)}
\end{split}
\end{equation*}
Inside each executor, PyArrow streams raw memory chunks to a \textbf{Numba JIT-compiled C-kernel} that computes the recursive state transitions for both daily Van Wagner (1987) and hourly FWI2025 algorithms. The terminal state vector $\mathcal{S}_{c,T}$ is persisted to a lightweight checkpoint table, while the computed FWI time-series are stored in Snappy-compressed columnar Parquet tables.

\subsubsection{Stage 4: Stateful Incremental Streaming and Spatial Validation}
When near-real-time observations arrive via the DMC API or operational forecasts, \textbf{Spark Structured Streaming} ingests the incremental batch, loads $\mathcal{S}_{c,T}$ from checkpoint storage, and updates the hazard surface in sub-second latency using \texttt{mapGroupsWithState}. Concurrently, \textbf{Apache Sedona} performs a distributed spatial join between high-FWI cells ($FWI > 50$) and NASA FIRMS active fire pixels / CONAF historical ignition logs to evaluate predictive skill.

\subsection{Simulation Methodology and Containerized Cluster Design}
To achieve rigorous, reproducible evaluation within student computational constraints, simulations will execute on an isolated \textbf{Docker-based cluster environment}:
\begin{itemize}
    \item \textbf{Topology:} 1 Spark Master container and $N \in \{1, 2, 4, 8\}$ Spark Worker containers, each constrained via Docker \texttt{cgroups} to strict CPU and RAM limits (e.g., 1 vCPU, 2~GB RAM per worker).
    \item \textbf{Storage Layer:} A local \textbf{MinIO} container providing S3-compliant distributed object storage, enabling realistic measurement of network socket I/O throughput.
    \item \textbf{Baseline Cluster:} A mirrored containerized \textbf{Dask/Xarray + Zarr} cluster subjected to identical resource caps.
\end{itemize}

\subsection{Experimental Evaluation Plan}
The forthcoming research phase will measure four core systems metrics:
\begin{enumerate}
    \item \textbf{Strong Scaling and Speedup:} Measuring total runtime $T(N)$ across $N \in \{1, 2, 4, 8\}$ workers for fixed decadal spans, evaluating empirical Speedup $S(N) = T(1)/T(N)$ against Amdahl's Law.
    \item \textbf{Shuffle Traffic and Read Amplification:} Monitoring Spark metrics (\texttt{ShuffleWriteBytes}, \texttt{SpillToDiskBytes}) to verify that spatial-first Parquet eliminates shuffle compared to naive time partitioning.
    \item \textbf{Storage and I/O Benchmarking:} Comparing disk footprint and random query latency across Uncompressed NetCDF, Flat CSV, Snappy Parquet, and ZSTD Parquet.
    \item \textbf{Domain Calibration:} Assessing Receiver Operating Characteristic (ROC-AUC) curves of calculated FWI against CONAF ignition coordinates during the 2017 and 2023 disaster windows.
\end{enumerate}

\section{Conclusion}
This paper establishes the state-of-the-art and technical solution architecture for distributed fire weather danger estimation across Chile. By revealing the data shortfalls of CR2MET v2.5 and formalizing the non-commutative temporal recurrence of the Canadian FWI, we demonstrated that the core computational bottleneck lies in the distributed corner-turn shuffle and multi-source spatial fusion. Our proposed Apache Spark and Parquet architecture with spatial-first partitioning eliminates network shuffling during recursive evaluation, providing a scalable, scientifically rigorous baseline for national wildfire resilience.

\balance
\bibliographystyle{IEEEtran}
\begin{thebibliography}{00}

\bibitem{cordero2024extreme}
R.~R. Cordero, S.~C. Feron, A.~Damiani, J.~F. Carrasco, C.~Karas, C.~Wang, C.~T. Kraamwinkel, and A.~Beaulieu, ``Extreme fire weather in Chile driven by climate change and El Ni\~no--Southern Oscillation (ENSO),'' \textit{Scientific Reports}, vol.~14, no.~1, p.~1974, 2024. [DOI: 10.1038/s41598-024-52481-x]

\bibitem{bowman2019human}
D.~M. J.~S. Bowman, A.~Moreira-Mu\~noz, C.~A. Kolden, R.~O. Ch\'avez, A.~A. Mu\~noz, F.~Salinas, A.~Gonz\'alez-Reyes, R.~Rocco, F.~de~la~Barrera, P.~M.~Freitas, and M.~E. Gonz\'alez, ``Human--environmental drivers and impacts of the globally extreme 2017 Chilean fires,'' \textit{Ambio}, vol.~48, no.~4, pp.~350--362, 2019. [DOI: 10.1007/s13280-018-1084-1]

\bibitem{garreaud2020central}
R.~D. Garreaud, J.~P. Boisier, R.~Rondanelli, A.~Montecinos, and H.~H. Sep\'ulveda, ``The Central Chile Mega Drought (2010--2018): A climate dynamics perspective,'' \textit{International Journal of Climatology}, vol.~40, no.~1, pp.~421--439, 2020. [DOI: 10.1002/joc.6219]

\bibitem{vanwagner1987development}
C.~E. Van~Wagner, ``Development and structure of the Canadian Forest Fire Weather Index System,'' Canadian Forestry Service, Petawawa National Forestry Institute, Chalk River, Ontario, Forestry Tech. Rep. 35, 1987.

\bibitem{wotton2009interpreting}
B.~M. Wotton, ``Interpreting and using outputs from the Canadian Forest Fire Danger Rating System in research applications,'' \textit{Environmental and Ecological Statistics}, vol.~16, no.~2, pp.~107--131, 2009. [DOI: 10.1007/s10651-007-0084-2]

\bibitem{boisier2023cr2met}
J.~P. Boisier, ``CR2MET: A high-resolution precipitation and temperature dataset for the period 1960--2021 in continental Chile (v2.5),'' Center for Climate and Resilience Research (CR2), Zenodo Data Set, 2023. [DOI: 10.5281/zenodo.7529682]

\bibitem{munoz2021era5land}
J.~Mu\~noz-Sabater, E.~Dutra, A.~Agust\'i-Panareda, C.~Albergel, G.~Arduini, G.~Balsamo, S.~Boussetta, M.~Choulga, S.~Harrigan, H.~Hersbach, \textit{et al.}, ``ERA5-Land: a state-of-the-art global reanalysis dataset for land applications,'' \textit{Earth System Science Data}, vol.~13, no.~9, pp.~4349--4383, 2021. [DOI: 10.5194/essd-13-4349-2021]

\bibitem{quilcaille2023fire}
Y.~Quilcaille, F.~Batibeniz, A.~F.~S. Ribeiro, M.~Hauser, and S.~I. Seneviratne, ``Fire weather index data under historical and shared socioeconomic pathway projections in the 6th phase of the Coupled Model Intercomparison Project from 1850 to 2100,'' \textit{Earth System Science Data}, vol.~15, no.~5, pp.~2153--2177, 2023. [DOI: 10.5194/essd-15-2153-2023]

\bibitem{dean2008mapreduce}
J.~Dean and S.~Ghemawat, ``MapReduce: simplified data processing on large clusters,'' \textit{Communications of the ACM}, vol.~51, no.~1, pp.~107--113, 2008. [DOI: 10.1145/1327452.1327492]

\bibitem{zaharia2012resilient}
M.~Zaharia, M.~Chowdhury, T.~Das, A.~Dave, J.~Ma, M.~McCauley, M.~J. Franklin, S.~Shenker, and I.~Stoica, ``Resilient distributed datasets: a fault-tolerant abstraction for in-memory cluster computing,'' in \textit{Proc. 9th USENIX Conference on Networked Systems Design and Implementation (NSDI '12)}, San Jose, CA, 2012, pp.~15--28.

\bibitem{armbrust2015sparksql}
M.~Armbrust, R.~S. Xin, C.~Lian, Y.~Huai, D.~Liu, J.~K. Bradley, X.~Meng, T.~Kaftan, M.~J. Franklin, A.~Ghodsi, and M.~Zaharia, ``Spark SQL: Relational data processing in Spark,'' in \textit{Proc. 2015 ACM SIGMOD International Conference on Management of Data}, Melbourne, Australia, 2015, pp.~1383--1394. [DOI: 10.1145/2723372.2742797]

\bibitem{armbrust2018structured}
M.~Armbrust, T.~Das, J.~Torres, B.~Yavuz, S.~Zhu, R.~Ramesh, A.~Narayan, K.~Mukherjee, B.~Joseph, A.~Rong, \textit{et al.}, ``Structured Streaming: A declarative API for real-time applications in Apache Spark,'' in \textit{Proc. 2018 ACM SIGMOD International Conference on Management of Data}, Houston, TX, 2018, pp.~601--613. [DOI: 10.1145/3183713.3190664]

\bibitem{carbone2017state}
P.~Carbone, S.~Ewen, G.~F\'ora, S.~Haridi, S.~Richter, and V.~Tzoumas, ``State management in Apache Flink: Consistent stateful distributed stream processing,'' \textit{Proceedings of the VLDB Endowment}, vol.~10, no.~12, pp.~1718--1729, 2017. [DOI: 10.14778/3137765.3137777]

\bibitem{melnik2010dremel}
S.~Melnik, A.~Gubarev, J.~J. Long, G.~Romer, S.~Shivakumar, M.~Tolke, and T.~Vassilakis, ``Dremel: Interactive analysis of web-scale datasets,'' \textit{Proceedings of the VLDB Endowment}, vol.~3, no.~1--2, pp.~330--339, 2010. [DOI: 10.14778/1920841.1920886]

\bibitem{palamuttam2015scispark}
R.~Palamuttam, H.~Mogrovejo, C.~A. Mattmann, B.~Wilson, P.~Ramirez, J.~Hartley, B.~Bue, S.~McGovern, and A.~J. Whitehall, ``SciSpark: Applying in-memory distributed computing to weather event detection and tracking,'' in \textit{Proc. 2015 IEEE International Conference on Big Data (Big Data)}, Santa Clara, CA, 2015, pp.~2020--2026. [DOI: 10.1109/BigData.2015.7363983]

\bibitem{hu2018climatespark}
F.~Hu, C.~Yang, D.~W. Schnase, and G.~D. Duffy, ``ClimateSpark: An in-memory distributed computing framework for big climate data analytics,'' \textit{Computers \& Geosciences}, vol.~115, pp.~111--121, 2018. [DOI: 10.1016/j.cageo.2018.03.011]

\bibitem{li2016spatiotemporal}
Z.~Li, F.~Hu, J.~L. Schnase, and G.~D. Duffy, ``A spatiotemporal indexing approach for efficient processing of big array-based climate data with MapReduce,'' \textit{International Journal of Geographical Information Science}, vol.~31, no.~1, pp.~17--35, 2016. [DOI: 10.1080/13658816.2015.1131830]

\bibitem{yu2018spatial}
J.~Yu, Z.~Zhang, and M.~Sarwat, ``Spatial data management in Apache Spark: The GeoSpark perspective and beyond,'' \textit{GeoInformatica}, vol.~23, no.~1, pp.~37--78, 2018. [DOI: 10.1007/s10707-018-0330-9]

\bibitem{gorelick2017google}
N.~Gorelick, M.~Hancher, M.~Dixon, S.~Ilyushchenko, D.~Thau, and R.~Moore, ``Google Earth Engine: Planetary-scale geospatial analysis for everyone,'' \textit{Remote Sensing of Environment}, vol.~202, pp.~18--27, 2017. [DOI: 10.1016/j.rse.2017.06.031]

\bibitem{rocklin2015dask}
M.~Rocklin, ``Dask: Parallel computation with blocked algorithms and task scheduling,'' in \textit{Proc. 14th Python in Science Conference (SciPy 2015)}, Austin, TX, 2015, pp.~130--136. [DOI: 10.25080/Majora-7b98e3ed-013]

\bibitem{hoyer2017xarray}
S.~Hoyer and J.~Hamman, ``xarray: N-D labeled arrays and datasets in Python,'' \textit{Journal of Open Research Software}, vol.~5, no.~1, p.~10, 2017. [DOI: 10.5334/jors.148]

\bibitem{abernathey2021cloud}
R.~P. Abernathey, T.~Augspurger, A.~Banihirwe, C.~Blackmon-Luca, T.~J. Crone, K.~Duffy, C.~Hamman, J.~Hamman, M.~J. Lepore, M.~R. Signell, \textit{et al.}, ``Cloud-native repositories for big scientific data,'' \textit{Computing in Science \& Engineering}, vol.~23, no.~2, pp.~26--35, 2021. [DOI: 10.1109/MCSE.2021.3059437]

\bibitem{cfs2025fwi}
Canadian Forest Service Fire Danger Group, ``The 2025 Update to the FWI System (FWI2025): Structure, changes and interpretation,'' Natural Resources Canada, Great Lakes Forestry Centre, Sault Ste. Marie, ON, Information Report GLC-X-42, 2025.

\bibitem{vitolo2020era5}
C.~Vitolo, F.~Di~Giuseppe, B.~Krzeminski, and J.~San-Miguel-Ayanz, ``ERA5-based global meteorological wildfire danger maps,'' \textit{Scientific Data}, vol.~7, no.~1, p.~216, 2020. [DOI: 10.1038/s41597-020-0554-z]

\bibitem{paperinstructions}
A.~Pizarro, E.~Valenzuela, and B.~Z\'u\~niga, ``Rubric and evaluation protocol for scientific literature triage in big data systems,'' Technical Specification, UDP-EIT-BD-2026-01, 2026.

\end{thebibliography}

\end{document}
